\section{Introduction and review question}

Contrastive vision--language models learn image and text representations that
can be compared in a shared embedding space. CLIP demonstrated the broad
transferability of this approach when both towers are jointly trained on
web-scale paired data \cite{radford2021clip}. Later work improved the objective,
data, distillation strategy, and deployment efficiency
\cite{zhai2023siglip,vasu2024mobileclip,faghri2025mobileclip2}. These gains,
however, are tied to substantial pretraining resources and tightly coupled
vision--language optimisation.

Independently pretrained unimodal encoders create a different opportunity. A
self-supervised vision model and a language model may already contain useful
semantic structure; if small trainable modules can align them, much of the cost
of joint pretraining can be avoided. This leads to the review question:

\begin{quote}
\emph{How far can independently pretrained frozen vision and language encoders
be aligned using a sub-five-million-parameter adaptation budget while
approaching compact jointly pretrained vision--language models under a fixed
inference budget?}
\end{quote}

The wording requires three distinct comparison groups. First, compact jointly
pretrained models such as MobileCLIP2-S0 are the most direct performance and
scale reference. Second, CLIP ViT-B/32 remains a useful field anchor because it
is widely recognised, but it should not be described as compact. Third,
Freeze-Align, SAIL, ShareLock, STRUCTURE, and SOTAlign form the methodological
peer group because they explicitly align frozen unimodal models
\cite{maniparambil2025freezealign,zhang2025sail,ruthardt2026sharelock,groger2025structure,roschmann2026sotalign}.
Conflating these groups can produce favourable but scientifically misleading
comparisons.

\section{Review method and scope}

This is a structured narrative review organised by technological decomposition:
problem, solution families, evaluation, and limitations. The search programme
was conducted in four rounds between 26 July and 10 August 2026. It combined a
ranked local bibliography, targeted OpenAlex searches, citation-graph screening,
and verification against primary paper or proceedings pages. The final refresh
used 20 decomposed academic-index searches, logged 322 records, and produced 273
globally unique records in its targeted export. It was considered alongside 54
ranked local entries and a separate screen of 288 works citing the closest
memory-bank papers.

Sources were prioritised when they were direct methodological peers, established
foundational work, or recent papers capable of changing the novelty boundary.
Published claims and numerical comparisons were checked against primary papers
or supplements before inclusion. Search summaries alone were not treated as
evidence. The review is comprehensive for the dissertation's defined questions,
but it is not a formal systematic review or meta-analysis: computer-science
indexing is incomplete, protocols are heterogeneous, and absence from the
screened corpus is not proof that no related work exists. The literature cut-off
is 10 August 2026.

\section{Jointly pretrained vision--language models}

\subsection{CLIP and sigmoid-objective models}

CLIP established the dominant dual-encoder formulation: an image encoder and a
text encoder are jointly trained so that matched pairs have high similarity and
other in-batch pairs act as negatives \cite{radford2021clip}. Its importance for
this project is twofold. It supplies the field's common retrieval vocabulary,
and OpenCLIP implementations make a recognisable ViT-B/32 reference available.
Nevertheless, CLIP B/32 is a field anchor rather than the compact model named by
the review question.

SigLIP replaces globally normalised softmax contrast with a pairwise sigmoid
loss, reducing dependence on a global similarity matrix and showing that gains
from ever-larger batches saturate \cite{zhai2023siglip}. SigLIP2 broadens this
recipe through captioning, self-supervised losses, online curation, multilingual
data, localisation, and multiple resolutions \cite{tschannen2025siglip2}. It is
therefore a useful alternative-objective reference, but its much larger B/32
stack is not a compact peer.

\subsection{Compact joint training and distillation}

Compact VLM research shows that parameter count alone is an incomplete measure
of difficulty. TinyCLIP combines affinity mimicking with weight inheritance
\cite{wu2023tinyclip}, while CLIP-KD evaluates several ways to transfer CLIP
knowledge into smaller models \cite{yang2024clipkd}. MobileCLIP introduces
multi-modal reinforced training and dataset reinforcement to improve mobile
image--text encoders \cite{vasu2024mobileclip}; MobileCLIP2 improves the same
training strategy and provides several compact reference scales
\cite{faghri2025mobileclip2}.

This literature has an important consequence for comparison. MobileCLIP2-S0 is
the primary compact jointly pretrained reference, not merely another teacher.
When it also supplies distillation targets to a student, it should be described
as both the teacher and an upper bound on what that selected teacher can
transfer. A smaller frozen-alignment system may still be scientifically useful,
but a claim of compact-model parity requires split-matched evaluation rather
than division of a validation score by a published test score.

\section{Frozen and locked alignment}

LiT showed that locking a pretrained image tower while training a text tower can
retain strong transfer performance \cite{zhai2022lit}. More recent work moves
closer to the present problem by starting with two independently pretrained
unimodal encoders. Freeze-Align learns lightweight mappings around frozen
encoders and studies pair compatibility \cite{maniparambil2025freezealign}.
SAIL also freezes the towers, but combines trainable alignment layers with a
refined sigmoid loss, multiple captions per image, pre-encoded features, and
very large effective batches \cite{zhang2025sail}. ShareLock emphasises the role
of strong language representations and reusable frozen features
\cite{ruthardt2026sharelock}. STRUCTURE adds geometry preservation when paired
data are limited \cite{groger2025structure}, while SOTAlign uses limited paired
data and larger unpaired collections through an optimal-transport objective
\cite{roschmann2026sotalign}.

These papers establish that frozen encoders plus small alignment heads are
prior art. Projectors, contrastive learning, compatibility screening, and
precomputed features are design choices rather than architectural
contributions. The appropriate research question is instead which design and
training conditions determine whether frozen alignment succeeds.
Table~\ref{tab:peers} summarises how each peer relates to the present setting.

\begin{table}[htbp]
\centering
\small
\caption{Direct frozen-alignment peers and their role in the comparison.}
\label{tab:peers}
\begin{tabularx}{\textwidth}{L{2.6cm}L{2.2cm}L{2.6cm}Y}
\toprule
\textbf{Work} & \textbf{Frozen components} & \textbf{Alignment regime} & \textbf{Relevance and limitation} \\
\midrule
Freeze-Align & Independently pretrained towers & Lightweight projectors; up to roughly 20M pairs & Direct compatibility and retrieval peer; larger encoders and data than the present constrained setting. \\
SAIL & Vision and language towers & Refined sigmoid alignment; CC12M; very large effective batch & Strong direct peer showing the benefit of objective and batch design. \\
ShareLock & Pretrained backbones & Lightweight alignment using precomputed features & Direct peer; published numbers depend on backbone and evaluation protocol. \\
STRUCTURE & Unimodal encoders & Geometry-preserving limited-data alignment & Direct conceptual peer; no single final Flickr configuration is sufficiently commensurate for scalar ranking. \\
SOTAlign & Unimodal encoders & Semi-supervised paired and unpaired alignment & Extends the data regime; recent forward-looking peer rather than a protocol-matched baseline. \\
\bottomrule
\end{tabularx}
\end{table}

Published Flickr30k retrieval figures illustrate the field but do not form a
common leaderboard. Freeze-Align reports image-to-text/text-to-image R@1 of
87.5/74.1 for its large final configuration
\cite{maniparambil2025freezealign}. In the SAIL supplement's common
DINOv2-B/CC12M comparison, SAIL-B reports 74.2/60.6 and ShareLock 68.1/49.3
\cite{zhang2025sail}. These systems differ in encoders, paired-data volume,
caption multiplicity, batch size, objectives, and split handling. Their numbers
answer ``where does frozen alignment research currently sit?'', not ``which row
wins under identical conditions?''

\section{Objectives, negatives, and historical embeddings}

\subsection{Why memory banks are attractive}

The contrastive objective used throughout this literature descends from InfoNCE,
which learns representations by discriminating a positive against a set of
sampled negatives \cite{oord2018cpc}. Its behaviour therefore depends directly
on how many negatives are available and how representative they are.

Contrastive alignment benefits from diverse negatives, but batch size is limited
by accelerator memory. Momentum Contrast introduced a queue-backed dictionary
whose key encoder changes slowly, making stored embeddings more consistent
\cite{he2020moco}. Cross-Batch Memory (XBM) generalised the use of historical
embeddings in supervised metric learning and argued that feature drift is slow
enough for recent stored representations to approximate current ones
\cite{wang2020xbm}. The method includes memory-size, batch-size, and drift
analyses, so neither queue-capacity curves nor the observation that embeddings
become stale is new.

Later work also demonstrates the limits of the slow-drift argument. Adaptive
Cross Batch Normalization treats current--memory mismatch as a distribution
alignment problem and reports conditions where memory performs worse than no
memory \cite{ajanthan2023acbn}. In cross-modal video--text retrieval, Memory
Enhanced Embedding Learning uses a momentum encoder precisely to prevent stored
features from evolving too rapidly \cite{zhao2021meel}. These precedents show
that stale memory can be harmful and that stabilisation mechanisms matter.

\subsection{The unresolved cross-modal question}

The closest prior work does not determine whether drift magnitude predicts
retrieval damage separately for image-side and text-side memories. XBM is
supervised and unimodal; its historical entries do not cross independently
parameterised modality-specific projectors. Cross-modal memory approaches use
memory as a component but generally evaluate final Recall@K rather than
separating measured age, representation drift, and degradation by modality.

Two works are the closest structural matches and were examined at full text.
CODER maintains dual dynamic modality-specific dictionaries for image--text
retrieval, so it possesses the apparatus required to ask the present question,
but it reports Recall@K and component-level ablations without a per-modality
bank ablation or any drift-versus-degradation analysis \cite{wang2022coder}.
CSMCIR explicitly identifies memory-bank staleness with age and responds with an
entropy-based dynamic bank, but it addresses composed image retrieval and again
does not measure drift as a quantity separate from degradation
\cite{qian2026csmcir}. Both therefore confirm that the ingredients are available
and that the specific decomposition has not been performed.

This yields the clearest gap in the reviewed corpus:

\begin{quote}
Prior work establishes that memory-bank embeddings drift and may become
harmful, but does not test whether equal measured drift implies equal damage
across the two towers of a frozen bidirectional image--text model.
\end{quote}

The distinction is important. ``Features move slowly'' is a statement about
representation change; ``stored negatives remain safe'' is a statement about
their effect on gradients and ranking. The first does not logically guarantee
the second, especially when the modalities occupy different geometries and the
alignment heads continue to learn.

\subsection{False negatives}

Image--caption relevance is many-to-many, while common benchmarks provide
non-exhaustive annotations. PCME models cross-modal examples probabilistically
and explicitly motivates its approach using ambiguity and incomplete relevance
labels \cite{chun2021pcme}. FALCON more directly balances hardness against
false-negative probability in vision--language contrastive learning
\cite{kim2026falcon}. This literature prevents a simplistic conclusion that
harder or more numerous negatives are always better.

For queue analysis, two false-negative mechanisms must be separated. Repeated
captions or images with a known shared image identity can be protected by a
multi-positive mask. Semantically related examples carrying different image
identifiers remain unlabelled and cannot be ruled out without semantic or human
relevance judgements. Consequently, a queue experiment can establish a
drift--damage dissociation without identifying a complete causal mechanism.

\section{Representation compatibility and the modality gap}

Frozen alignment assumes that independently trained encoders expose compatible
semantic structure. The Platonic-representation hypothesis motivates this possibility by arguing
that models trained on different data and objectives converge towards a shared
statistical model of reality \cite{huh2024platonic},
while frozen-alignment studies test it empirically through encoder-pair selection
and geometry measures. Freeze-Align, for example, uses representational
similarity to reason about pair quality \cite{maniparambil2025freezealign}.
However, compatibility should not be treated as an intrinsic property of an
encoder pair: rankings may also depend on the alignment objective, caption
handling, batch construction, and optimisation recipe.

The modality-gap literature provides a second constraint. Image and text
embeddings can occupy separated regions or narrow cones even after contrastive
training, with initialisation and optimisation contributing to the separation
\cite{liang2022gap}. More recent work provides a gradient-flow account and
shows that mismatched pairs and temperature can influence the gap
\cite{yaras2025modalitygap}. Therefore equal cosine movement need not have equal
meaning in the two towers.

The properties of the frozen vision substrate are a third constraint, and they
are fixed before alignment begins. Self-supervised families such as DINOv2 and
DINOv3 are trained without language supervision and are the encoders that make
frozen alignment attractive in the first place
\cite{oquab2023dinov2,simeoni2025dinov3}. Because their weights are not updated,
choices made during their pretraining---patch size, positional encoding
mechanism, and training resolution---continue to govern what a downstream
alignment head can read. An input resolution that departs from the encoder's
pretraining regime can therefore change retrieval quality for reasons that have
nothing to do with the alignment module, which makes resolution and
preprocessing policy part of the experimental description rather than an
implementation detail.

This is a rival explanation, not a reason to abandon the queue question. A
static difference in geometry predicts a baseline or scale difference; it does
not by itself predict how retrieval degradation changes with queue age. A
defensible analysis must therefore compare response rates over overlapping
drift ranges and describe the result as modality-specific sensitivity, not as a
fully identified new mechanism.

\section{Reading information from frozen tokens}

Global encoder outputs are convenient but may discard token-level information.
Set Transformer introduced pooling by multi-head attention using learned seed
vectors \cite{lee2019settransformer}. CoCa uses attentional poolers within a
joint image--text model \cite{yu2022coca}; BLIP-2's Q-Former reads a frozen image
encoder through learned queries \cite{li2023blip2}; and SigLIP2 uses MAP pooling
\cite{tschannen2025siglip2}. Attentive probing examines lightweight attention
readers over frozen image features explicitly as an accuracy--parameter trade-off
\cite{psomas2026attention}.

Learned-query aggregation is consequently an established primitive. A study
that compares CLS or mean pooling against learned-query or transformer readers
can still make a useful empirical contribution: it can establish whether useful
information is present in frozen patch or word tokens and how much trainable
capacity is required to expose it. It should not claim that a learned query, MAP
head, or Q-Former-like reader is a new architecture.

The distinction between access and selection is especially useful. If mean
pooling produces little improvement but learned interaction produces a much
larger gain, the evidence suggests that token-level information exists but must
be selected or combined conditionally. That result concerns the behaviour of
the frozen representation under a budget, not ownership of the pooling block.

\section{Parameter-efficient adaptation}

Strict freezing tests whether fixed representations already contain alignable
information, but it may impose an avoidable ceiling. LoRA established the
now-standard alternative: freeze the pretrained weights and train injected
low-rank update matrices, which decouples adaptation capacity from the size of
the host model \cite{hu2021lora}. VL-Adapter demonstrates adapter-based
parameter-efficient transfer for vision--language tasks
\cite{sung2022vladapter}. More recent work moves closer to the present dual-tower
setting. HALoRA allocates low-rank budgets hierarchically while accounting for
vision--language asymmetry \cite{zhang2026halora}. B-HFA combines shared adapters
with hierarchical visual aggregation for retrieval, although it adapts an
already pretrained VLM rather than aligning two independent unimodal towers
\cite{xu2026bhfa}.

Low-rank adaptation of one or both towers is therefore not algorithmically new.
Its value in a constrained dissertation is as a controlled upper bound: a
vision-only, text-only, and dual-tower factorial can quantify where strict
freezing costs performance and whether mismatch is isolated to either modality.
The strictly frozen endpoint and the low-rank endpoint must remain separate
scientific categories, and trainable parameters must be reported separately
from the total deployed stack.

\section{Token reduction and deployment efficiency}

Reducing vision-token count is a mature efficiency strategy. DynamicViT learns
to prune tokens progressively \cite{rao2021dynamicvit}; ToMe merges similar
tokens without retraining \cite{bolya2022tome}; Token Fusion combines pruning
and merging according to model sensitivity \cite{kim2023tokenfusion}; and PuMer
directly prunes and merges tokens in vision--language models
\cite{cao2023pumer}. Recent parameter-free or saliency-driven work includes
Dyna-ViT and SAD-TM \cite{rubab2026dynavit,xie2026sadtm}.

This literature establishes two requirements for an efficiency claim. First,
lower FLOPs or fewer tokens do not guarantee lower wall-clock latency because
kernel shape, memory movement, padding, and batch scheduling also matter.
Second, a faster model is not on a better efficiency frontier if retrieval
accuracy falls disproportionately. Token reduction should therefore be measured
end to end and compared with a matched no-reduction parent on both accuracy and
latency.

The present study supplies direct evidence for both requirements, reported in
full in Sections 4.9, 4.10, and 5.6. A fixed $2\times2$ internal merge reduced
third-quartile latency by 0.98--1.61\,ms, confirming that internal token count
is latency-relevant, but a four-arm three-seed evaluation against matched
no-merge parents cost between 5.69 and 15.95 percentage points of retrieval
accuracy. The accuracy loss therefore dwarfs the latency gain at every merge
depth and under both the strictly frozen and low-rank parents, and token
reduction did not improve the measured accuracy--latency frontier. The fixed
merge is accordingly a simple baseline against the adaptive literature rather
than a new pruning method, and the adaptive-routing branch remains unrun.

\section{Data scale, curation, and evaluation comparability}

Reproducible scaling studies show power-law relationships among model size,
paired data, and compute, while also showing that the training distribution
changes those relationships \cite{cherti2023scaling}. DataComp demonstrates
that curation strategies can substantially alter quality under controlled
training and compute \cite{gadre2023datacomp}. DataComp-VLM extends the
data-centric perspective to broader mixtures and downstream task suites
\cite{farina2026datacompvlm}. A small alignment model trained on COCO
\cite{lin2014coco} should not be expected to reproduce all capability acquired
by a jointly trained web-scale model, and a remaining reference gap cannot be
attributed solely to architecture.

The evaluation surface is similarly narrow. Flickr30k \cite{young2014flickr30k}
and COCO supply short descriptive captions, so bidirectional R@1 on these sets
measures short-caption retrieval and nothing more. Zero-shot classification,
compositional reasoning, and domain shift are distinct capabilities that a
retrieval score does not certify, and results established on one should not be
transferred to the others by implication. A study confined to retrieval is
legitimate provided the resulting claims are stated at that scope.

Evaluation protocol is equally consequential. At least five quantities must be
kept distinct:

\begin{enumerate}
  \item validation performance used during model selection;
  \item sealed or final-test performance;
  \item published external results using different data and splits;
  \item locally trainable inference parameters versus total deployed parameters;
  \item latency measured in the same allocation versus historical timing from a
        different session or software stack.
\end{enumerate}

A ratio between validation performance and an external test score is not a
split-matched retention measure. Similarly, ``one third the parameters'' is
meaningless unless the numerator and denominator count the same deployed
components. Strong efficiency evidence reports hardware, precision, batch size,
input and padding policy, warm-up, repetitions, measurement order, and paired
uncertainty.

\section{Data governance, inherited bias, and licensing}

Efficiency framing can obscure a question that the reuse of pretrained
components makes unavoidable: what is inherited along with the weights. A frozen
encoder carries the representational consequences of its pretraining corpus, and
because those weights are never updated, alignment cannot correct them. Any bias
present in the vision or language tower is therefore transmitted into the shared
space rather than diluted by it, and a small trained projector is a poor place to
look for a remedy.

The data-centric literature reinforces this. Curation decisions are shown to
change downstream quality substantially \cite{gadre2023datacomp}, which is also
to say that they determine what a model has and has not seen. Web-scale corpora
are assembled by automated filtering rather than editorial review, so the
resulting coverage is uneven in ways that a retrieval score will not surface.
Reporting accuracy while remaining silent on provenance describes only half of
what a system does.

Three practical constraints follow for work of this kind. First, performance and
safe deployment are separate claims: a favourable Recall@K supports the former
and says nothing about the latter. Second, pretrained checkpoints and datasets
arrive under licences that govern redistribution, so the defensible position is
to reference upstream artefacts and release code and configurations rather than
mirror weights or images. Third, an honest account distinguishes what was
measured from what was merely not examined; this review does not constitute a
demographic fairness audit, and no such audit should be inferred from the
efficiency and retrieval evidence assembled here.

\section{Synthesis: what is established and what remains open}

Table~\ref{tab:synthesis} separates what the reviewed corpus already settles
from the questions that remain open for this dissertation.

\begin{longtable}{L{2.7cm}L{4.5cm}L{5.4cm}}
\caption{Synthesis of the literature and the remaining research opportunity.}
\label{tab:synthesis}\\
\toprule
\textbf{Area} & \textbf{Established by prior work} & \textbf{Status in this dissertation} \\
\midrule
\endfirsthead
\toprule
\textbf{Area} & \textbf{Established by prior work} & \textbf{Status in this dissertation} \\
\midrule
\endhead
Frozen alignment & Frozen towers can be connected with lightweight projectors; compatibility and limited-data objectives matter. & Resolved (Sections 4.2, 4.7, 5.1): a strictly frozen endpoint was built and its recipe diagnosed under the stated budget. \\
Historical memory & Queues enlarge the negative set; drift, distribution mismatch, and memory harm are known. & Resolved and central (Sections 4.3, 5.2): equal measured drift did not imply equal damage, and the two towers differed sharply in sensitivity. \\
Modality geometry & Image and text embeddings can occupy different regions of the shared space. & Partly resolved (Sections 4.3, 5.2): an overlap-restricted slope comparison supports modality-specific sensitivity, but geometry-normalised drift remains future work and the mechanism is not identified. \\
Token aggregation & Learned-query and attentive pooling are established. & Resolved (Sections 4.6, 5.3): the internal ablation quantifies what a small reader recovers; no architectural novelty is claimed. \\
Distillation & Compact CLIP distillation is mature. & Resolved as a bounded result (Section 4.5): teacher choice was tested at fixed small-data alignment; tuned mixtures and ensembles were not searched. \\
Low-rank adaptation & Adapters and LoRA are established for VLM transfer. & Resolved (Sections 4.8, 5.4): on the sealed test the bounded low-rank relaxation gains 9.30pp over the strictly frozen endpoint, quantifying the cost of strict freezing. \\
Token reduction & Adaptive pruning and merging can improve efficiency. & Resolved negatively (Sections 4.9, 4.10, 5.6): merging saved 0.98--1.61\,ms but cost 5.69--15.95pp, so the frontier did not improve; adaptive routing remains unrun. \\
Data and measurement & Scale, distribution, curation, and deployment protocol materially change outcomes. & Resolved (Sections 4.5, 4.11, 5.5): a sealed split-matched test and corrected same-allocation profiling now separate what survives from what does not. \\
Task scope & Retrieval benchmarks use short descriptive captions and do not certify classification, compositional, or shifted-domain ability. & Open by disclosure (Section 6.2): zero-shot classification diagnostics exist but the evaluation package explicitly declines to treat them as confirmatory tests. \\
Governance & Curation shapes coverage; frozen weights transmit inherited bias; licences constrain redistribution. & Open by design: no fairness audit, licence opinion, or deployment-readiness claim is made. \\
\bottomrule
\end{longtable}

The most defensible unoccupied contribution is the memory-queue question. A
generic statement that stale embeddings hurt would duplicate XBM and ACBN. The
narrower contribution is to test the sufficiency of slow drift as a safety
argument in a bidirectional frozen cross-modal regime: measure realised queue
age, representation drift, and retrieval degradation independently, then test
whether the relationship differs by modality. This is a falsifiable boundary
condition on established memory-bank reasoning rather than a claim that all
queues are harmful.

The second contribution is empirical rather than architectural. Comparing
global summaries, mean pooling, learned queries, and transformer aggregation
under the same frozen towers can reveal whether alignment fails because
information is absent or merely difficult to read. The relevant novelty is the
controlled finding and its budgeted frontier, not the attention-pooling module.

The remaining experiments---data coverage, distillation, low-rank adaptation,
latency repair, and token reduction---complete the explanation by testing rival
hypotheses and establishing boundaries. Their value lies in a coherent chain of
evidence in which negative results constrain the final claim.

\section{Position of the dissertation}

The literature supports a methodological and empirical dissertation, not an
architectural novelty claim. The project is relevant because it studies the
space between two mature but imperfectly connected research lines: large-scale
joint VLM pretraining and lightweight alignment of existing unimodal models.
Its central question is practically important for researchers who cannot repeat
web-scale pretraining, while its queue analysis addresses a more general
scientific issue: whether a commonly monitored proxy, representation drift, is
sufficient to certify the safety of historical negatives.

The comparison hierarchy should remain explicit:

\begin{enumerate}
  \item \textbf{Primary compact reference:} MobileCLIP2-S0, also disclosed as
        the distillation teacher.
  \item \textbf{Field anchor:} OpenCLIP ViT-B/32, useful for recognisability and
        local same-allocation comparisons but not labelled compact.
  \item \textbf{Alternative objective:} SigLIP2 ViT-B/32.
  \item \textbf{Methodological peers:} Freeze-Align, SAIL, ShareLock,
        STRUCTURE, and SOTAlign.
\end{enumerate}

Within that hierarchy, a claim that a model is approximately one third the
size is defensible only against OpenCLIP B/32 when both values count the full
deployed stacks. It is not defensible against MobileCLIP2-S0. Likewise, an
accuracy-retention ratio is defensible only for a split-matched local comparison.
These restrictions make the resulting contribution more credible: the thesis
does not claim that a small model beats CLIP, but explains how much alignment is
recoverable, which interventions supply it, which efficiency shortcuts fail,
and where the remaining gap comes from.

\section{Limitations of the literature evidence}

The evidence base has six limitations. First, frozen-alignment papers use
different encoders, paired datasets, caption multiplicities, objectives, and
evaluation splits, preventing a formal meta-analysis. Second, recent 2026 work
may not yet be fully indexed. Third, computer-science proceedings are unevenly
covered by citation indexes, so the absence of a matching experiment is a
bounded corpus statement. Fourth, published parameter counts do not always
separate trainable and total inference parameters. Fifth, accuracy, FLOPs,
latency, training energy, and deployment memory are not interchangeable notions
of efficiency. Sixth, the reviewed evidence is overwhelmingly retrieval-centred
and reports little on inherited bias or licensing, so conclusions about
deployment suitability cannot be drawn from it.

These limitations lead to conservative language. The review can establish that
frozen alignment and learned aggregation are prior art, that queue staleness is
known, and that the per-modality drift--damage decomposition is not present in
the screened corpus. It cannot prove universal novelty, compact-reference
parity, or deployment readiness. The search should be refreshed immediately
before final dissertation submission.

\section{Conclusion}

Efficient frozen vision--language alignment is relevant not because it replaces
joint pretraining outright, but because it tests what can be recovered from
existing unimodal representations under realistic adaptation constraints. The
literature establishes the components---projectors, contrastive objectives,
learned queries, distillation, LoRA, and token reduction---and therefore narrows
the dissertation's contribution to how those components behave when controlled
carefully. The strongest research gap concerns historical embeddings: prior work
uses slow drift to justify memory, but does not show that equal drift produces
equal consequences across modalities. A staged experimental study of that gap,
followed by controlled aggregation, adaptation, data, and latency analyses,
provides a coherent and appropriately bounded contribution to efficient
vision--language learning.
